\newpage

\subsection{Data}
\label{subsec:data}
AI applications stand and fall with the training data. Despite that, there has only recently been a paradigm shift towards the datasets due to the data-centric AI perspective. The AI Act requires in Article 10 (3) that \emph{training, validation, and testing datasets shall be relevant, representative, free of errors and complete}. What this means in practice, however, is often unclear. Our extension of a data card builds on previous work on data documentation but is extended with further references, especially those addressing regulation, whenever possible \citep{Pushkarna2022_DataCards, Gebru2021_DataSheets}. From the standardization side, the \cite{ISO5259} series deals with data quality for AI systems, \cite{ISO8183} with the data life cycle framework for AI systems and \cite{ISO25012} describes a general data quality model. The first two are still under development but are worth reviewing once they are final. Before introducing the data card, we discuss open questions that we think are not covered sufficiently in the state of the art. 


\subsubsection{Related Research} 
We start by summarizing research directions that might be worth considering before starting data collecting.
Two related active areas of research are \emph{active learning} \citep{Settles2010_ActiveLearningSurvey} and \emph{core sets} \citep{Feldman2020_CoreSetsSurvey}. In active learning, the algorithm is given access to a small labeled dataset $D$ and a large unlabeled dataset $U$. The algorithm can then query an oracle (e.g., human labeler) for labels up to some budget $b$. The goal is to optimize model performance within the  limit of the budget. A core set is a small fraction $S\in D$ of the original dataset $D$ such that the learning algorithm achieves a similar performance on this subset as if it was trained on the entire dataset.

Recently, the concept of memorization or uniqueness of data points was introduced \citep{Feldman2020_LongTail, Jiang2020_CharacterizingRegularities}.  A data instance $i$ from the training set is called unique if removing it from the training set reduces the probability for $i$ to be classified correctly by the model. \cite{Feldman2020_LongTail} hypothesizes that such rare points are important for the generalization of models. I.e., a single rare point such as an image of a red car could determine whether this concept is correctly classified by the model in the real world or not. 
\citet{Jiang2020_CharacterizingRegularities} use it as a measure to categorize the structure of a dataset and show that mislabeled points are harder to memorize. Unfortunately, estimating these instances is computationally expensive. For each data instance at least two models need to be trained; one with and the other without the instance. However, there might be occasions when it could become useful, for example, if a specific point shall be inspected. 
Finally, \citet{Meding2021_DDD} describe dichotomous data-difficulty and show that many datasets such as ImageNet suffer from imbalanced data difficulty. There are many data points in the test set that are never classified correctly (called \emph{impossible}) and many that are always classified correctly (called \emph{trivial}). They show that models can better be compared on the remaining points.  A similar conclusion can be drawn for label errors. \cite{Northcutt2021_LabelErrors} show that larger models tend to be favored on datasets with label errors while smaller models might actually outperform them when evaluated on a dataset without errors.

These research fields and recent advances are related to data collection and quality and performing a literature review before starting to collect data can help to make the collection more effective. 


%\pagebreak
\subsubsection{Data Collection}
\label{subsec:data_collection}
After a literature review, the first step is to gather data. Apart from collecting metadata and protected attributes, we focus on suggestions around dataset size and coverage. It is worth mentioning that, to the best of our knowledge, there are currently no established practices in that regard and the upcoming ISO standards on data \cite{ISO5259} will most likely not include detailed requirements. Hence, the following suggestions should be seen as points worth to consider rather than being mandatory. 

\paragraph{Metadata and Leakage} Although being well-understood, leakage remains a central problem in many datasets \citep{Kapoor22_LeakageAndReproducibility}. A famous example is CheXNet for pneumonia detection on x-ray images \citep{Rajpurkar2017_CheXNet}, where the authors performed a na\"ive train-test split with overlapping x-ray images of the same patients in training and test sets in the first version of the paper. As such, the performance of the model on real data was overestimated. 
When creating train and test splits, such leaks can be avoided by collecting metadata and using it for stratification or to generate data splits of different difficulty.  In the ChexNet example, the metadata could be an anonymous ID for each patient and an easy test set would contain data from patients already represented in the train data wheres a hard test set would only contain data of unseen patients. This information can help to estimate the generalization capabilities of a model which can differ between applications. If the goal is to predict re-occurrence of a disease, performance on the easy test set could be the primary metric while performance on the hard test set would be of interest if the goal is to detect pneumonia in new patients. In general, it is advisable to collect whatever metadata is easy to obtain and compatible with data protection law because it is often impossible to backtrack once the dataset is final. This especially includes protected attributes.  

\hspace*{10mm}\subparagraph{Protected Attributes}
\label{app:protected_attrs}
Usually, the protected attributes include sex, race, skin colour, ethnic or social origin, 
genetic features, language, religion or belief, political or any other opinion, membership of a national minority, property, birth, disability, age or sexual orientation and can be found, for instance, in Art. 21 of the Charter of Fundamental Rights of the European Union \citep{CharterFundamentalRightsEU}. According to the EU handbook on non-discrimination law, 
it might differ between use-cases whether an attribute is considered as protected \citep{Handbook2018_NonDiscrimination}. Furthermore, pregnancy can be considered as protected in some cases. Access to such information is crucial when detecting biases but there may also be cases when other laws prohibit the use of such attributes \citep{van_Bekkum_2023_DoesGDPRNeedException}. 

\paragraph{Dataset Size}
 One problem when collecting data is knowing when a sufficient amount was acquired. Despite the standard perspective that ``there's no data like more data'', according to IBM's Robert Mercer, collecting and annotating data is both expensive and not always feasible.  For example, in medicine it is almost impossible to increase the number of samples of rare diseases. Data augmentation is one solution, but in this section we want to focus on the size of the raw dataset.  A rule of thumb in computer vision suggests to use around 1.000 images per class for image classification tasks as a good starting point. The one-in-ten rule\footnote{\url{https://en.wikipedia.org/wiki/One_in_ten_rule}} suggests that per model-parameter ten data points are necessary. One way to estimate the necessary amount of data is by tracking model improvements when sequentially adding larger fractions of the training data. It is well known that the resulting performance plot behaves logarithmically \citep{Figueroa2012_PredictingSampleSize, Viering2021_LearningCurveShape}. Hence, the largest improvements are obtained early on. As shown in Figure \ref{fig:performance_train_fraction}(a), with about 50\% of the training data, the performance is already above 90\% and doubling the amount of data increases the performance by less than 5\%. In practice, one approach could be to start with a small test set for data collection and utilize it to estimate when a reasonable amount of data will be gathered. Once the data has been acquired this test set may be added to the train data or discarded and a new test set is created.

 \begin{figure}
    \centering
    \subfloat[{\scriptsize Model Performance for Train Fraction}]{{
        \includegraphics[width=0.45\textwidth]{train_size_perf.pdf}
    }}
    \subfloat[{\scriptsize Test Variance for Test Fraction}]{{
        \includegraphics[width=0.45\textwidth]{test_var.pdf}
    }}
    
    \caption{(a) Model performance of a \ac{MLP} trained on the MNIST dataset with increasing fraction of original train data. The improvement of adding data points behaves logarithmically. %\\ \\
    %\hspace*{9mm}
    (b) Variance over 20 runs of training an \ac{MLP} on the entire MNIST training set and evaluating the performance on fractions of the test set. With increasing size of the test set, the variance decreases.}
    \label{fig:performance_train_fraction}
\end{figure}

\paragraph{Data Coverage}
Apart from the size of the dataset, covering all possible scenarios is important. Recent research has shown that many data points are redundant   \citep{Toneva2018_catastrophicForgetting, Sorscher2022_BeyondScalingLaws, Feldman2020_CoreSetsSurvey} and hypothesised that unique points (which may be considered as edge cases) can have a significant impact on model performance  \citep{Feldman2020_LongTail}. One way to find such samples is by manually defining edge cases, e.g. under the involvement of domain experts, and actively collecting these data points. This can be extended by a combination of empirical and technical analysis, for instance, by plotting a downsampled embedding and analyzing the edge cases. A useful tool for this purpose is  \href{https://tensorleap.ai/}{tensorleap}. 
\pagebreak
Unfortunately, it is still unclear whether rare points are desirable or not in a dataset. While they can have a significant impact on model performance, they might make a model prone to privacy attacks \citep{Sorscher2022_BeyondScalingLaws, Feldman2020_LongTail}. 


%\pagebreak
\subsubsection{Train-Test Split Creation}
Usually, after data collection the data would be labled and preprocessed. Since our suggestion regrading labeling are short, they are covered in the section on documentation. Apart from this, we have swapped the section on train-test split creation and preprocessing because it is recommend to do preprocessing on the splits separately whenever feasible.

\paragraph{Stratification}
\label{app:stratification}
First of all, training and test sets should be stratified.
If metadata was collected it should be considered for stratification. As indicated in Section \ref{subsec:data_collection}, another option is to create splits of different difficulty according to the metadata. An easy test set might contain overlapping instances between entities in training and test set (e.g., data of the same patients in training and test set) whereas a hard test set contains only data from  new entities (patients). This can be used to estimate the generalization capabilities of a model and helps to understand whether it is usable for a task or not. 

\paragraph{Test Set Size}
One of the most difficult choices is the size of the test set.
Despite the standard 80-20 (or similar) suggestion we are not aware of common strategies to estimate the necessary test set size in a general setting \citep{Joseph2022_OptimalRatioSplitting}. In medical and psychological research, sample size estimation based on statistical significance is established. For example, if the effect of medication against a certain disease should be evaluated and it is a priori known that the disease occurs in 10 from 1000 patients ($p=1\,\%$) and the target is an error margin of $\epsilon=$1\,\% with a confidence level of $\sigma=$ 99\,\%, the necessary sample size $n$ is given by 
$n \approx \frac{z_{0.99}^2\times p(1-p)}{\epsilon^2}=\frac{2.58^2 \times 0.01(1-0.01)}{0.01^2}=659$, 
where $z_{0.99}=2.58$ is the value for the selected confidence interval derived from the normal distribution. If $p$ is unknown, it is usually set to 0.5.
% https://www.questionstar.de/blog/berechnung-der-stichprobengrose-technische-details-einfach-erklaert/#:~:text=der%20notwendige%20Stichprobenumfang.-,Beispiel%3A,%2C05%5E2%20%3D%20385.
The same can be applied to \ac{ML} and was actually discussed in a few works especially from the medical domain \citep{Beleites2013_SampleSizePlanning, Dietterich_ApproximateStatisticalTestsSupervisedClassification, Raudys1990_SmallSampleSizeReccomendationPractitioners, Konietschke2021_SmallSampleSize}. 
In the case of ML, the interpretation is as follows; the goal is to evaluate an ML model that predicts the disease with a randomly drawn test set of size 659. Then, there is a 99\,\% confidence that the real world model performance is $\pm$1\,\% of the error rate on the test set. However, one drawback is that it is assumed that the collected data is representative of the real world setting and it is not guaranteed that the model will not learn any spurious correlations. 

Another indicator for a good test set is the test variance. The underlying assumption is that a better test set should show a low variance between different evaluation runs with the same model \citep{Bouthillier2021_Variance}. That is, if the same experiment results in an unexpectedly large variance, the test set is most likely too small and should be extended. An example for this effect is shown in Figure \ref{fig:performance_train_fraction}(b). Although multiple evaluations should generally be done, it is of special importance if the test set cannot be extended \citep{Bouthillier2021_Variance, Raschka2018_ModelEvaluationAndSelection, Lones2021_AvoidMLPitfalls}.


\subsubsection{Data Processing}
\label{subsec:data_processing}
After the data was collected and (maybe) labeled, it is preprocessed. This, for example, includes data cleaning, feature selection, and feature engineering. As mentioned above, it is important to do this separately on all splits whenever possible \citep{Lones2021_AvoidMLPitfalls, poretschkin2021_IAIS_Leitfaden}. Otherwise, information can leak into the test set. For example, by oversampling rare classes before creating splits \citep{winter2021_TueV_Austria}.  
In this work, we distinguish between processing steps that are a fixed part of a deployed dataset and as such is documented in the data card,
and steps specific to a particular \ac{ML} model which are documented in the model card. Generally, it is advisable to involve domain experts and to document whether and how the processing steps are agreed on with them \citep{poretschkin2021_IAIS_Leitfaden}. 


\subsubsection{Understanding and Documenting Data \& Label Quality}
\label{subsec:data_quality}
Finally, the characteristics of the data need to be documented.
This includes all measures applied in order to understand the data. For example, plots of embeddings, identified rare points, or possible label errors or ambiguities. Although not focused on AI, the 16 dimensions defined in the \cite{ISO25012} standard are well suited for this. They include accuracy, completeness, consistency, credibility, currentness, accessibility, compliance, confidentiality, efficiency, precision, tractability, understandability, availability, portability, recoverability, and relevance. Dimensions that are of particular importance are accuracy (can be interpreted as fraction of label errors), completeness (fraction of missing values) and consistency (fraction of inconsistencies such as synonyms).
 
From the technical side, \href{https://cleanlab.ai/}{Cleanlab} \citep{Northcutt2021_LabelErrors} and \href{https://tensorleap.ai/}{tensorleap} are toolboxes that provide various methods for analyzing data and detecting and fixing label errors. Although it is not always possible to estimate whether something is a label error, documenting such ambiguous samples is important. Apart from this, the four-eyes principle is advisable when data is annotated manually \citep{poretschkin2021_IAIS_Leitfaden}. 
When dealing with third-party annotators, the usage of a \emph{gold standard dataset} is possible.
A gold-standard dataset is a small fraction of data with very high label quality that is send to an external provider for labeling and can be used to estimate the label quality of the provider afterwards. 

\newpage
{\small

\renewcommand\arraystretch{1.3}
\begin{longtable}{p{0.15\linewidth}|p{0.35\linewidth}p{0.15\linewidth}p{0.35\linewidth}}  
    \toprule
    \multicolumn{4}{c}{\bf{Data Card}}\\
    \toprule
    \bf{Step}        & \bf{Requirement}                                         &\bf{AI Act}       & \bf{References} \\
    \hline
    \hline
    \multirow{5}{=}{General}
                        & Name originator of the dataset and provide a contact person          & & \cite{Pushkarna2022_DataCards, Gebru2021_DataSheets}\\\cline{2-4}
                        & Describe the intended use of the dataset                             & & \cite{Pushkarna2022_DataCards, Gebru2021_DataSheets}\\\cline{2-4}
                        & Describe licensing and terms of usage                                & & \cite{Pushkarna2022_DataCards, Gebru2021_DataSheets} \\\cline{2-4}
                        
    \specialrule{.075em}{0em}{0em}
    \multirow{16}{=}{Data-description}
                        & Collect requirements for the data before starting data collection  & Art. 10 (2) d) e) \\\cline{2-4}
                        & Describe a data point with its interpretation                      & & \cite{Pushkarna2022_DataCards, Gebru2021_DataSheets} \\\cline{2-4}
                        & Maybe, provide additional documentation to understand the data (e.g., links to scientific sources, preprocessing steps or other necessary information) & & \cite{Pushkarna2022_DataCards}   \\\cline{2-4}
                        & If there is GDPR relevant data (i.e. personally identifiable information), describe it  & & \cite{Pushkarna2022_DataCards, Gebru2021_DataSheets}\\\cline{2-4}
                        & If there is biometric data, describe it                       &                   & \cite{Pushkarna2022_DataCards, Gebru2021_DataSheets, poretschkin2021_IAIS_Leitfaden}\\\cline{2-4}
                        & If there is copyrighted data, summarize it                    & Art. 28 b)        & \\\cline{2-4}
                        & If there is business relevant information, describe it        &                   & \cite{poretschkin2021_IAIS_Leitfaden}\\\cline{2-4}

    \specialrule{.075em}{0em}{0em}
    \multirow{16}{=}{Collection}
                        & Describe the data collection procedure and the data sources   & Art. 10 (2) b)  Annex III c)  & \cite{Pushkarna2022_DataCards, Gebru2021_DataSheets}\\\cline{2-4}
                        & Use data version control                                      & Art. 10 (2) b)    & \cite{poretschkin2021_IAIS_Leitfaden} \\\cline{2-4}
                        & Consider the prior requirements for the data                  & Art. 10 (2) e)    & \\\cline{2-4}
                        & Include and describe metadata                                 &                   & \href{https://www.dcc.ac.uk/guidance/standards/metadata}{Metadata Standards}\\\cline{2-4}
                        & Involve domain experts and describe their involvement         &                   & \cite{poretschkin2021_IAIS_Leitfaden} \\\cline{2-4}
                        & Describe technical measures to ensure completeness of data& Art.  10 (2) g) Art. 28 b) & \href{https://github.com/cleanlab/cleanlab/}{CleanLab}\\\cline{2-4}
                        & Describe and record edge cases                                &                   &  \cite{poretschkin2021_IAIS_Leitfaden, Pushkarna2022_DataCards} \\\cline{2-4}
                        & If personal data is used, make sure and document that all individuals know they are part of the data         & & \cite{Lu2021_DataDeprecation, Gebru2021_DataSheets} \\\cline{2-4}
                        & Describe if and how the data could be misused                 &                   & \cite{Gebru2021_DataSheets, Pushkarna2022_DataCards}\\\cline{2-4}
                        \pagebreak
                        \cline{2-4}
                        & If fairness is identified as a risk, list sensitive attributes & Art. 10 (5)  & Section \ref{app:protected_attrs}, \cite{ISO12791_Fairness} \\\cline{2-4}
                        & If applicable, address data poisoning                          &              & \cite{poretschkin2021_IAIS_Leitfaden}\\

    \specialrule{.075em}{0em}{0em}
    \multirow{4}{=}{Labeling}
                        & If applicable, describe the labeling process                   & & \cite{Gebru2021_DataSheets, Pushkarna2022_DataCards}, \href{https://snorkel.ai/}{snorkel.ai}\\\cline{2-4}
                        & If applicable, describe how the label quality is checked       & & Section \ref{subsec:data_quality}\\
                        
    \specialrule{.075em}{0em}{0em}
    \multirow{13}{=}{Splitting}
                        & Create and document meaningful splits with stratification      & Art.  10 (3)      &  Section \ref{app:stratification},  \cite{Beleites2013_SampleSizePlanning, Konietschke2021_SmallSampleSize, Raudys1990_SmallSampleSizeReccomendationPractitioners, Joseph2022_OptimalRatioSplitting, Zhang2022_SliceTeller, poretschkin2021_IAIS_Leitfaden} \\\cline{2-4}
                        & Describe how data leakage is prevented                                  &                   & \cite{poretschkin2021_IAIS_Leitfaden, winter2021_TueV_Austria}\\\cline{2-4}
                        & Recommendation: test the splits and variance via cross-validation       &                   & \cite{Bouthillier2021_Variance}\\\cline{2-4}
                        & Recommendation: split dataset into difficult, trivial and moderate      &                   & \cite{Meding2021_DDD} \\\cline{2-4}
                        & Recommendation: put special focus on label quality of test data         &                   & \cite{Northcutt2021_LabelErrors}\\\cline{2-4}
                        & Reminder: Perform separate data preprocessing on the splits             &                   & \cite{Roberts2017_CrossValidationTemporalSpatial, Lones2021_AvoidMLPitfalls}\\
                        
    \specialrule{.075em}{0em}{0em}
    \multirow{8}{=}{Preprocessing}
                        & Document and motivate all processing steps that are a fixed part of the data      & Art. 10 (2) c)    & Section \ref{subsec:data_processing} \\\cline{2-4}
                        & Document whether the raw data can be accessed                                     &                   & \cite{Gebru2021_DataSheets}\\\cline{2-4}
                        & If sensitive data is available highlight (pseudo)-anonymization                   & Art.  10 (5)      & \cite{poretschkin2021_IAIS_Leitfaden}\\\cline{2-4}
                        & If fairness is a risk, highlight fairness specific preprocessing                  & Art.  10 (5)      & \cite{poretschkin2021_IAIS_Leitfaden}\\
                        
    \specialrule{.075em}{0em}{0em}
    \multirow{9}{=}{Analyzing}
                        & Understand and document characteristics of test and training data.                                            &   Art. 10 (2) e) g), Art 10 (3) & Section \ref{subsec:data_quality}, \cite{ISO25012}, \cite{ISO5259}   \cite{Zhang2022_SliceTeller, Mitchell2022_MeasuringData, Mazumder2022_DataPerf} \\\cline{2-4} 
                        & Document why the data distribution fits the real conditions or why this is not necessary for the use case     & Art. 10 (4) & \cite{poretschkin2021_IAIS_Leitfaden}\\\cline{2-4}
                        & Document limitations such as errors, noise, bias or known confounders   & Art. 10 (2) f) g) & \cite{Gebru2021_DataSheets, Pushkarna2022_DataCards, Wittenbrink2022_LeitfadenQualitatsmanagementKI, poretschkin2021_IAIS_Leitfaden} \\

    \specialrule{.075em}{0em}{0em}
    \pagebreak
    \specialrule{.075em}{0em}{0em}
    \multirow{10}{=}{Serving}
                        & Describe how the dataset will be maintained in future                &               & \cite{Gebru2021_DataSheets, Pushkarna2022_DataCards}, \cite{Lu2021_DataDeprecation}\\\cline{2-4}
                        & Describe the storage concept (e.g., everything users need to know to access the data). For developers it must be possible to document on which version of the data a specific model was trained.                                 &               & \cite{poretschkin2021_IAIS_Leitfaden}\\\cline{2-4}
                        & Describe the backup procedure                                        &               & \cite{poretschkin2021_IAIS_Leitfaden}\\\cline{2-4}
                        & If necessary, document measures against data poisoning               &               & \cite{poretschkin2021_IAIS_Leitfaden}\\

    \specialrule{.075em}{0em}{0em}
    \multirow{2}{=}{Further notes}       
                        & Document further recommendations or shortcomings in the data \\\cline{2-4}
                        
    \specialrule{.075em}{0em}{0em}
    \multirow{2}{=}{Involvement of Individuals}                            
                        & If feasible, describe involvement of affected individuals            & Art. 29 5a. \\     

    \bottomrule
  \end{longtable}
}

